% --- Archon physics formalization source begin ---
% archon:physics
% archon:covers PhyXMiniProblems/problem_phyx_mini_0074.lean
% archon:source-report reports/phyx_mini/problem_phyx_mini_0074.source.json

\chapter{Physics problem phyx\_mini\_0074}
\label{ch:PhyXMiniProblems_problem_phyx_mini_0074}

\paragraph{Problem source.}
\#\# Physical scenario

There’s one angle of incidence \textbackslash{}beta onto a prism for which the light inside an isosceles prism travels parallel to the base and emerges at angle \textbackslash{}beta.\textbackslash{} A laboratory measurement finds that \textbackslash{}beta = 52.2\textasciicircum{}\textbackslash{}circ for a prism shaped like an equilateral triangle.

\#\# Question

What is the prism’s index of refraction?

\#\# Answer choices

- A: A: 1.47

- B: B: 2.26

- C: C: 1.58

- D: D: 1.95

\#\# Figure caption (auxiliary; use the image as primary evidence)

The image shows a transparent triangle with three important angles labeled. 

- A triangle (likely a prism) is shaded light blue.
- The apex angle of the triangle at the top is labeled \textbackslash{}(\textbackslash{}alpha\textbackslash{}).
- The triangle also has a ray entering from the left and exiting to the right, bending inside the triangle.
- The angles of incidence and refraction at the entry and exit points are labeled \textbackslash{}(\textbackslash{}beta\textbackslash{}).
- A normal line is drawn inside the triangle, perpendicular to the sides at the entry and exit points.
- The text "n" appears inside the triangle, likely representing the refractive index of the prism material.

\#\# Dataset metadata

Geometrical Optics; Physical Model Grounding Reasoning, Spatial Relation Reasoning

\paragraph{Recorded answer/context.}
C: C: 1.58

\paragraph{Figure/image path.}
/root/proposal\_for\_physic/science-mango/phyx\_mini\_run/phyx\_data/test\_image/74.png

\paragraph{Formalization target.}
create a compiling Lean file with sorry bodies at `PhyXMiniProblems/problem\_phyx\_mini\_0074.lean`. The Lean declarations must preserve the physical quantities, dimensions or dimensional roles, figure labels, governing-law hypotheses, and final relation expressed by this problem.
Use Mathlib/Physlib names found through LeanExplore where available. If a physics API is missing, introduce faithful local abstractions rather than scalar placeholder aliases.

\begin{theorem}[Physics formalization target]
\label{thm:physics:phyx_mini_0074:target}
\lean{PhyXMiniProblems.ProblemPhyxMini0074.prismRefractiveIndexIsAnswerC}
\uses{def:physics:phyx-mini-0074:phyxminiproblems-problemphyxmini0074-anglefromdegrees, def:physics:phyx-mini-0074:phyxminiproblems-problemphyxmini0074-equilateralprismdiagram, def:physics:phyx-mini-0074:phyxminiproblems-problemphyxmini0074-matchesprimaryfigure, def:physics:phyx-mini-0074:phyxminiproblems-problemphyxmini0074-hasphysicalopticalparameters, def:physics:phyx-mini-0074:phyxminiproblems-problemphyxmini0074-matchesparallelbaseraygeometry, def:physics:phyx-mini-0074:phyxminiproblems-problemphyxmini0074-obeyssnellslaw, def:physics:phyx-mini-0074:phyxminiproblems-problemphyxmini0074-answerchoice, def:physics:phyx-mini-0074:phyxminiproblems-problemphyxmini0074-roundstonearesthundredth, def:physics:phyx-mini-0074:phyxminiproblems-problemphyxmini0074-isnearestanswerchoice}
The assigned autoformalize agent should translate this physics problem into Lean declarations in the covered file, with theorem and lemma proof bodies written as `by sorry`.
\end{theorem}
\begin{proof}
This is an autoformalization task, not a proof task. Produce statements that can later be proved by the physics prover without weakening the physical model.
\end{proof}

% --- Archon declaration topology begin ---
\section{Declaration topology}
\begin{definition}[angle From Degrees]
  \label{def:physics:phyx-mini-0074:phyxminiproblems-problemphyxmini0074-anglefromdegrees}
  \lean{PhyXMiniProblems.ProblemPhyxMini0074.angleFromDegrees}
  Convert a degree readout to Mathlib's angle type
\end{definition}
\begin{proof}
  This is the declared model definition of angle From Degrees.
\end{proof}
\begin{definition}[refractive Index Dimension]
  \label{def:physics:phyx-mini-0074:phyxminiproblems-problemphyxmini0074-refractiveindexdimension}
  \lean{PhyXMiniProblems.ProblemPhyxMini0074.refractiveIndexDimension}
  Refractive index is dimensionless in Physlib's dimension system
\end{definition}
\begin{proof}
  This is the declared model definition of refractive Index Dimension.
\end{proof}
\begin{definition}[Diagram Plane]
  \label{def:physics:phyx-mini-0074:phyxminiproblems-problemphyxmini0074-diagramplane}
  \lean{PhyXMiniProblems.ProblemPhyxMini0074.DiagramPlane}
  The Cartesian plane containing the triangular prism cross-section
\end{definition}
\begin{proof}
  This is the declared model definition of Diagram Plane.
\end{proof}
\begin{definition}[Ray Segment]
  \label{def:physics:phyx-mini-0074:phyxminiproblems-problemphyxmini0074-raysegment}
  \lean{PhyXMiniProblems.ProblemPhyxMini0074.RaySegment}
  The three directed portions of the depicted light ray
\end{definition}
\begin{proof}
  This is the declared model definition of Ray Segment.
\end{proof}
\begin{definition}[Prism Face]
  \label{def:physics:phyx-mini-0074:phyxminiproblems-problemphyxmini0074-prismface}
  \lean{PhyXMiniProblems.ProblemPhyxMini0074.PrismFace}
  The three faces of the equilateral triangular cross-section
\end{definition}
\begin{proof}
  This is the declared model definition of Prism Face.
\end{proof}
\begin{definition}[Refracting Interface]
  \label{def:physics:phyx-mini-0074:phyxminiproblems-problemphyxmini0074-refractinginterface}
  \lean{PhyXMiniProblems.ProblemPhyxMini0074.RefractingInterface}
  The two refracting interfaces crossed by the ray
\end{definition}
\begin{proof}
  This is the declared model definition of Refracting Interface.
\end{proof}
\begin{definition}[Optical Medium]
  \label{def:physics:phyx-mini-0074:phyxminiproblems-problemphyxmini0074-opticalmedium}
  \lean{PhyXMiniProblems.ProblemPhyxMini0074.OpticalMedium}
  The homogeneous optical media traversed by the ray
\end{definition}
\begin{proof}
  This is the declared model definition of Optical Medium.
\end{proof}
\begin{definition}[Figure Angle Label]
  \label{def:physics:phyx-mini-0074:phyxminiproblems-problemphyxmini0074-figureanglelabel}
  \lean{PhyXMiniProblems.ProblemPhyxMini0074.FigureAngleLabel}
  The three angle labels visible in the primary image
\end{definition}
\begin{proof}
  This is the declared model definition of Figure Angle Label.
\end{proof}
\begin{definition}[incident Medium]
  \label{def:physics:phyx-mini-0074:phyxminiproblems-problemphyxmini0074-incidentmedium}
  \lean{PhyXMiniProblems.ProblemPhyxMini0074.incidentMedium}
  \uses{def:physics:phyx-mini-0074:phyxminiproblems-problemphyxmini0074-refractinginterface, def:physics:phyx-mini-0074:phyxminiproblems-problemphyxmini0074-opticalmedium}
  Incident medium along the propagation direction at each interface
\end{definition}
\begin{proof}
  This is the declared model definition of incident Medium.
\end{proof}
\begin{definition}[transmitted Medium]
  \label{def:physics:phyx-mini-0074:phyxminiproblems-problemphyxmini0074-transmittedmedium}
  \lean{PhyXMiniProblems.ProblemPhyxMini0074.transmittedMedium}
  \uses{def:physics:phyx-mini-0074:phyxminiproblems-problemphyxmini0074-refractinginterface, def:physics:phyx-mini-0074:phyxminiproblems-problemphyxmini0074-opticalmedium}
  Transmitted medium along the propagation direction at each interface
\end{definition}
\begin{proof}
  This is the declared model definition of transmitted Medium.
\end{proof}
\begin{definition}[Equilateral Prism Diagram]
  \label{def:physics:phyx-mini-0074:phyxminiproblems-problemphyxmini0074-equilateralprismdiagram}
  \lean{PhyXMiniProblems.ProblemPhyxMini0074.EquilateralPrismDiagram}
  \uses{def:physics:phyx-mini-0074:phyxminiproblems-problemphyxmini0074-diagramplane, def:physics:phyx-mini-0074:phyxminiproblems-problemphyxmini0074-raysegment, def:physics:phyx-mini-0074:phyxminiproblems-problemphyxmini0074-prismface, def:physics:phyx-mini-0074:phyxminiproblems-problemphyxmini0074-refractinginterface, def:physics:phyx-mini-0074:phyxminiproblems-problemphyxmini0074-opticalmedium, def:physics:phyx-mini-0074:phyxminiproblems-problemphyxmini0074-figureanglelabel}
  Physical quantities and labelled geometry carried by the prism diagram. At each interface, incidence and refraction angles are measured from the depicted normal on the incident and transmitted sides respectively. A Module.Ray records a nonzero direction modulo positive rescaling, so equality of the internal-ray and base directions faithfully expresses the oriented parallelism shown in the figure
\end{definition}
\begin{proof}
  This is the declared model definition of Equilateral Prism Diagram.
\end{proof}
\begin{definition}[Matches Primary Figure]
  \label{def:physics:phyx-mini-0074:phyxminiproblems-problemphyxmini0074-matchesprimaryfigure}
  \lean{PhyXMiniProblems.ProblemPhyxMini0074.MatchesPrimaryFigure}
  \uses{def:physics:phyx-mini-0074:phyxminiproblems-problemphyxmini0074-anglefromdegrees, def:physics:phyx-mini-0074:phyxminiproblems-problemphyxmini0074-equilateralprismdiagram}
  Primary-image and laboratory readouts. These premises record the equilateral prism, the meanings of α and the two β labels, the measured 52.2° external angles, and the air-index calibration. They do not assign a value to the prism's refractive index
\end{definition}
\begin{proof}
  This is the declared model definition of Matches Primary Figure.
\end{proof}
\begin{definition}[Is Physical Ray Angle]
  \label{def:physics:phyx-mini-0074:phyxminiproblems-problemphyxmini0074-isphysicalrayangle}
  \lean{PhyXMiniProblems.ProblemPhyxMini0074.IsPhysicalRayAngle}
  An incidence or refraction angle is on the acute physical branch shown in the ray diagram
\end{definition}
\begin{proof}
  This is the declared model definition of Is Physical Ray Angle.
\end{proof}
\begin{definition}[Has Physical Optical Parameters]
  \label{def:physics:phyx-mini-0074:phyxminiproblems-problemphyxmini0074-hasphysicalopticalparameters}
  \lean{PhyXMiniProblems.ProblemPhyxMini0074.HasPhysicalOpticalParameters}
  \uses{def:physics:phyx-mini-0074:phyxminiproblems-problemphyxmini0074-equilateralprismdiagram, def:physics:phyx-mini-0074:phyxminiproblems-problemphyxmini0074-isphysicalrayangle}
  Positivity and branch conditions for the elementary optical model
\end{definition}
\begin{proof}
  This is the declared model definition of Has Physical Optical Parameters.
\end{proof}
\begin{definition}[Matches Parallel Base Ray Geometry]
  \label{def:physics:phyx-mini-0074:phyxminiproblems-problemphyxmini0074-matchesparallelbaseraygeometry}
  \lean{PhyXMiniProblems.ProblemPhyxMini0074.MatchesParallelBaseRayGeometry}
  \uses{def:physics:phyx-mini-0074:phyxminiproblems-problemphyxmini0074-equilateralprismdiagram}
  Geometric relations induced by the internal ray being parallel to the base of an equilateral prism. The two internal normal angles are equal, and the standard prism-angle relation says that their sum is the apex angle. No field fixes a refractive index or selects an answer choice
\end{definition}
\begin{proof}
  This is the declared model definition of Matches Parallel Base Ray Geometry.
\end{proof}
\begin{theorem}[internal Interface Angles Are Thirty Degrees]
  \label{thm:physics:phyx-mini-0074:phyxminiproblems-problemphyxmini0074-internalinterfaceanglesarethirtydegrees}
  \lean{PhyXMiniProblems.ProblemPhyxMini0074.internalInterfaceAnglesAreThirtyDegrees}
  \uses{def:physics:phyx-mini-0074:phyxminiproblems-problemphyxmini0074-anglefromdegrees, def:physics:phyx-mini-0074:phyxminiproblems-problemphyxmini0074-equilateralprismdiagram, def:physics:phyx-mini-0074:phyxminiproblems-problemphyxmini0074-matchesprimaryfigure, def:physics:phyx-mini-0074:phyxminiproblems-problemphyxmini0074-hasphysicalopticalparameters, def:physics:phyx-mini-0074:phyxminiproblems-problemphyxmini0074-matchesparallelbaseraygeometry}
  For the acute physical branch, equal internal angles whose sum is the 60° apex angle are both 30°. This isolates the figure-derived geometry from the subsequent application of Snell's law
\end{theorem}
\begin{proof}
  The conclusion follows from the governing relations and figure readouts named in the statement of internal Interface Angles Are Thirty Degrees.
\end{proof}
\begin{definition}[Snells Law At]
  \label{def:physics:phyx-mini-0074:phyxminiproblems-problemphyxmini0074-snellslawat}
  \lean{PhyXMiniProblems.ProblemPhyxMini0074.SnellsLawAt}
  \uses{def:physics:phyx-mini-0074:phyxminiproblems-problemphyxmini0074-refractinginterface, def:physics:phyx-mini-0074:phyxminiproblems-problemphyxmini0074-incidentmedium, def:physics:phyx-mini-0074:phyxminiproblems-problemphyxmini0074-transmittedmedium, def:physics:phyx-mini-0074:phyxminiproblems-problemphyxmini0074-equilateralprismdiagram}
  Snell's law n₁ sin θ₁ = n₂ sin θ₂ at one labelled interface
\end{definition}
\begin{proof}
  This is the declared model definition of Snells Law At.
\end{proof}
\begin{definition}[Obeys Snells Law]
  \label{def:physics:phyx-mini-0074:phyxminiproblems-problemphyxmini0074-obeyssnellslaw}
  \lean{PhyXMiniProblems.ProblemPhyxMini0074.ObeysSnellsLaw}
  \uses{def:physics:phyx-mini-0074:phyxminiproblems-problemphyxmini0074-equilateralprismdiagram, def:physics:phyx-mini-0074:phyxminiproblems-problemphyxmini0074-snellslawat}
  The depicted light ray obeys Snell's law at both prism faces
\end{definition}
\begin{proof}
  This is the declared model definition of Obeys Snells Law.
\end{proof}
\begin{definition}[Answer Choice]
  \label{def:physics:phyx-mini-0074:phyxminiproblems-problemphyxmini0074-answerchoice}
  \lean{PhyXMiniProblems.ProblemPhyxMini0074.AnswerChoice}
  Labels of the four displayed answer choices
\end{definition}
\begin{proof}
  This is the declared model definition of Answer Choice.
\end{proof}
\begin{definition}[refractive Index Readout]
  \label{def:physics:phyx-mini-0074:phyxminiproblems-problemphyxmini0074-answerchoice-refractiveindexreadout}
  \lean{PhyXMiniProblems.ProblemPhyxMini0074.AnswerChoice.refractiveIndexReadout}
  \uses{def:physics:phyx-mini-0074:phyxminiproblems-problemphyxmini0074-answerchoice}
  Dimensionless refractive-index readout printed for each answer choice
\end{definition}
\begin{proof}
  This is the declared model definition of refractive Index Readout.
\end{proof}
\begin{definition}[recorded Answer Choice]
  \label{def:physics:phyx-mini-0074:phyxminiproblems-problemphyxmini0074-recordedanswerchoice}
  \lean{PhyXMiniProblems.ProblemPhyxMini0074.recordedAnswerChoice}
  \uses{def:physics:phyx-mini-0074:phyxminiproblems-problemphyxmini0074-answerchoice}
  Dataset metadata: the recorded answer label is C. This is not used as a premise of the target theorem
\end{definition}
\begin{proof}
  This is the declared model definition of recorded Answer Choice.
\end{proof}
\begin{definition}[Rounds To Nearest Hundredth]
  \label{def:physics:phyx-mini-0074:phyxminiproblems-problemphyxmini0074-roundstonearesthundredth}
  \lean{PhyXMiniProblems.ProblemPhyxMini0074.RoundsToNearestHundredth}
  A physical index rounds to the displayed value to the nearest hundredth
\end{definition}
\begin{proof}
  This is the declared model definition of Rounds To Nearest Hundredth.
\end{proof}
\begin{definition}[Is Nearest Answer Choice]
  \label{def:physics:phyx-mini-0074:phyxminiproblems-problemphyxmini0074-isnearestanswerchoice}
  \lean{PhyXMiniProblems.ProblemPhyxMini0074.IsNearestAnswerChoice}
  \uses{def:physics:phyx-mini-0074:phyxminiproblems-problemphyxmini0074-answerchoice}
  A displayed choice is strictly nearer to the physical index than every other displayed choice
\end{definition}
\begin{proof}
  This is the declared model definition of Is Nearest Answer Choice.
\end{proof}
% --- Archon declaration topology end ---
% --- Archon physics formalization source end ---
